\chapter{DevOps Implementation}

\section{Dockerization Strategy}
The project is container-ready with separate services for MySQL, backend API, and frontend web server.

\begin{itemize}[leftmargin=1.3cm]
\item \texttt{backend/Dockerfile}: builds Node.js backend image and runs \texttt{server.js} on port 5000.
\item \texttt{frontend/Dockerfile}: builds React app and serves static output with Nginx.
\item \texttt{frontend/nginx.conf}: handles SPA fallback routing and proxies \texttt{/api} requests to backend service.
\end{itemize}

\section{Docker Compose Orchestration}
\texttt{docker-compose.yml} defines three services:
\begin{itemize}[leftmargin=1.3cm]
\item \textbf{mysql}: MySQL 8 container with initialization from \texttt{database/} scripts.
\item \textbf{backend}: Depends on healthy MySQL and exposes port 5000.
\item \textbf{frontend}: Depends on healthy backend and exposes port 80.
\end{itemize}

\section{Health Checks}
\begin{itemize}[leftmargin=1.3cm]
\item MySQL health check uses \texttt{mysqladmin ping}.
\item Backend health check calls \texttt{GET /api/health} from inside container.
\item Health endpoint runs \texttt{SELECT 1} query to confirm database connectivity.
\item If database check fails, endpoint returns HTTP 503 with \texttt{database: disconnected}.
\end{itemize}

\section{GitHub Actions CI Pipeline}
\texttt{.github/workflows/ci.yml} automatically runs on push and pull request.

\begin{itemize}[leftmargin=1.3cm]
\item Sets up Node.js 20.
\item Installs backend dependencies using \texttt{npm install}.
\item Verifies backend app import with \texttt{node -e "require('./backend/src/app')"}-style check.
\item Installs frontend dependencies and executes \texttt{npm run build}.
\end{itemize}

This pipeline quickly catches import errors and frontend build failures before deployment.

\section{Windows Automation Scripts}
\subsection*{1. One-click startup script}
\texttt{start-project.bat}:
\begin{itemize}[leftmargin=1.3cm]
\item Detects Node/npm availability.
\item Installs backend/frontend dependencies only if \texttt{node\_modules} is missing.
\item Starts backend and frontend in separate CMD windows.
\item Opens browser at \texttt{http://localhost:5173}.
\end{itemize}

\subsection*{2. Backup script}
\texttt{scripts/backup-db.bat}:
\begin{itemize}[leftmargin=1.3cm]
\item Uses full-path \texttt{mysqldump.exe}.
\item Creates \texttt{backups/} if missing.
\item Produces timestamped SQL dump file for \texttt{bharat\_finvest\_ops}.
\end{itemize}

\subsection*{3. Restore script}
\texttt{scripts/restore-db.bat}:
\begin{itemize}[leftmargin=1.3cm]
\item Accepts drag/drop backup file path from user.
\item Uses full-path \texttt{mysql.exe} to restore selected dump into target database.
\end{itemize}

\section{Environment-based Configuration}
Backend configuration uses environment variables from \texttt{backend/.env}:
\texttt{DB\_HOST}, \texttt{DB\_PORT}, \texttt{DB\_USER}, \texttt{DB\_PASSWORD}, \texttt{DB\_NAME}, \texttt{JWT\_SECRET}, \texttt{JWT\_EXPIRES\_IN}, and \texttt{PORT}.

This design supports local execution, Docker execution, and future production deployment with minimal changes.
